\begin{abstract}
\noindent
La fiscalité des revenus du capital est censée orienter l'épargne vers le financement
productif, mais on observe mal où va l'argent qu'une baisse d'impôt libère. Le prélèvement
forfaitaire unique instauré en France en 2018, une imposition proportionnelle des dividendes et
des plus-values à 30 \%, offre un cas net. Nous comparons les entreprises françaises à
celles de cinq économies comparables, avant et après la réforme, sur un panel de plus de
16 millions d'années-sociétés tiré d'Orbis, complété par les sociétés cotées de Compustat
Global et par les comptes nationaux d'Eurostat. La réforme n'a pas fait entrer de capital neuf
dans les fonds propres des sociétés, ni privées ni cotées, et n'a provoqué ni substitution par la
dette ni hausse de l'investissement. L'argent est sorti vers les ménages sous forme de dividendes,
en hausse de 32 \% l'année de la réforme, puis s'est logé en épargne liquide et en
portefeuille d'actions, pas dans le capital des entreprises ni dans l'immobilier. En plaçant côte à
côte la hausse d'impôt de 2013 et la baisse de 2018 dans un cadre unique, nous mettons au jour une
asymétrie. La hausse de 2013 avait réalloué du capital vers les bénéfices retenus ; la baisse
symétrique de 2018 n'en ramène qu'une fraction, moins de la moitié une fois neutralisée
l'anticipation. Le capital part plus facilement qu'il ne revient. Ce cliquet fiscal, cohérent sur
la distribution et sur l'allocation du capital, montre qu'une baisse d'impôt ne défait pas les
déformations d'une hausse passée, un résultat qui importe pour toute évaluation de la fiscalité du
capital.
\end{abstract}
